\section{Scope, evidence and outstanding tasks}\label{sec:tasks}

\subsection{Control families and the meaning of the results}

The complete Green identity is an exact representation; the stability
and residual floors are quantitative scoped obstructions; and
Corollary~\ref{cor:canonical-log} is an explicit local acquisition
estimate. These roles must not be interchanged. In particular a
lower bound for a sufficient quadratic residual is not a lower
bound for \(J_0\).

The collision controls summarized in \cite{OGControls} give the
following analytical screen. No new control computation is required
for these premise comparisons.

\begin{center}\small
\begin{tabularx}{\textwidth}{@{}>{\raggedright\arraybackslash}p{0.24\textwidth}Y@{}}
\toprule
Control family & Premise not supplied for the canonical application\\
\midrule
Euler-factor and translated-source families &
The extra multiplier changes the literal theta/xi field. Canonical
prime subtraction and Euler-anchor constants cannot simply be assigned
to the altered function. The fixed-prime baseline in
Section~\ref{sec:conditioning} is itself not theta.\\
Weighted-square, finite-prefix and finite-prime-orbit controls &
Matched finite data do not supply the complete normalized theta
coefficients together with full reflection.\\
Reflected super-Gaussian control &
Reflection and positive Gram forms do not supply the integer-square
theta expansion or the unchanged xi/zeta identity.\\
Smooth shape and positive-primitive controls &
Their specified shape, endpoint and asymptotic data do not supply the
complete canonical arithmetic source.\\
Atomic, finite-profile and variable-readout examples &
Only their actual stated source and observation premises may be used;
the full canonical source/profile/cosine-readout conjunction is absent.\\
\bottomrule
\end{tabularx}
\end{center}

The generic local theorem can apply to an admissible control with its
own anchor and growth constants. If that control has an off-axis zero
in the relevant domain, accurate acquisition retains its positive
Green signal and cannot pass the corresponding strict exclusion
threshold. Failure of the literal theta identity does not imply
failure of every inequality derived from it. A new canonical sign
estimate would have to identify its quantitative distinguishing
premise.

\subsection{Evidence and assistance}

This draft was prepared for Edward Baker with substantial OpenAI Codex
LLM assistance in derivation, exposition, source checking and production.
The current session's recorded history supplies no model or
reasoning-effort configuration; none is inferred. The source/resolvent
conditioning and short-mollifier transfer retain their internal
assessment with the additive correction
\cite{OGCorrection,OGAssessment}. The local boundary-log theorem
retains proof-author checking \cite{OGBoundaryLog}.
The parent Green representation and prime reduction retain their
existing scoped internal assessments
\cite{LocalizedMetricsAssessment,PrimeBaselineAssessment}.
The external moment theorem is an explicit imported input.
There has been no independent specialist human review or formal
verification. Manuscript integration does not strengthen these
evidence levels.

No numerical experiment is needed for any proof in this draft.
Future computations can support feasibility, error-budget evaluation
or a declared finite certificate, but cannot supply evidence for an
all-height RH-equivalent statement merely by checking finite cases.

\subsection{Outstanding analytical tasks}

\paragraph{1. Complete the direct logarithmic assembly.}
The existing frozen execution task \cite{OGAssembly} asks for a
finite-source upper approximation to \(J_0(X)\).
For integer \(R\ge1\), \(Q\ge2\), and
\(\sup_{|x-X|\le R}|F-\widehat F|\le\delta\), its proposed form is
\begin{align}\label{eq:proposed-U}
 U_{R,Q,\delta}(X)
 ={}&\int_{X-R}^{X+R}K(x-X)\log(|\widehat F(x)|+\delta)\dd x
 \nonumber\\
 &+\mathcal T(X,R)-P_Q(X)+\beta_Q+G(X,1).
\end{align}
Here \(\widehat F=D\widehat H\) uses the unchanged source,
and \(\mathcal T\) must bound the absolute omitted logarithmic tail
of the exact \(F\). The finite compact error sum proposed from
\eqref{eq:canonical-envelope} is
\[
 \mathcal E_{X,R}(\delta)=
 \sum_{k=0}^{R-1}
 \frac{\mathcal E_{X+k+1/2}(\delta)+
       \mathcal E_{X-k-1/2}(\delta)}{16\cosh(ak)}.
\]
The outstanding complete statement is
\begin{equation}\label{eq:proposed-assembly}
 0\le U_{R,Q,\delta}(X)-J_0(X)
 \le\mathcal E_{X,R}(\delta)+2\mathcal T(X,R)+2\beta_Q.
\end{equation}
Equations~\eqref{eq:proposed-U}--\eqref{eq:proposed-assembly}
are a specified \emph{target}, not a theorem proved by this manuscript.
The assembly must derive the full explicit log tail, verify the
compact multiplier/source payment, and justify the ordered finite
choices of \(R,Q,\delta,N_{\rm src},V\).
Both factors of two are intentional: replacing a signed tail by an
upper allowance can cost twice its absolute bound, and the same is
true for prime truncation. The actual signed compact integral must
be retained.

\paragraph{2. Assess the assembled claim and its local input.}
A claim-level assessment should recheck the line-sublevel theorem,
its canonical application, all source and logarithmic tails,
the single plus pole correction, and evaluation-error propagation.
The local result currently has proof-author checking; the complete
assembly has not yet been executed. Changing the model alone would
not supply specialist independence.

\paragraph{3. Acquire a favorable canonical value.}
Even after \eqref{eq:proposed-assembly}, one needs a directed evaluated
upper value satisfying
\[
 U\le\overline U\le U+\eta_{\rm eval},\qquad
 \overline U<m(d,r).
\]
Alternatively, a source-specific analytical estimate could establish
the same strict upper comparison directly.
An approximation error tending to zero does not force this value
inequality. The missing information is a quantitative bound on the
signed critical-line logarithmic average relative to its exact prime
baseline, including the exceptional centers relevant to a proposed
zero.

\paragraph{4. Resolve uniformity and practical precision.}
The present local exponent decreases with height; the compact
source multiplier in \eqref{eq:multiplier} is deliberately coarse.
One must separate finite acquisition for each prescribed tuple
from common estimates as \(|X|\to\infty\) or \(d\downarrow0\).
Sharpening constants may improve a practical diagnostic but does not
establish canonical smallness. Any new resolvent or contour route
must pay its denominators, residues, harmonic data and tails.

\paragraph{5. Use numerical evidence only for a declared finite question.}
After the analytical assembly and a finite error budget are fixed,
a bounded diagnostic may measure cost and certified accuracy against
a comparable direct method. It should have specified cases and a
stopping rule. No pilot, all-height extrapolation, favorable canonical
margin or new numerical claim is included in the present draft.

